\title{QLoRA: Efficient Finetuning of Quantized LLMs}

\begin{abstract}
We introduce \textsc{QLoRA}, a memory-efficient finetuning method that enables finetuning of a 65 billion parameter model on a single 48GB GPU, while maintaining the full performance of 16-bit finetuning tasks. \textsc{QLoRA} performs backpropagation through a frozen, 4-bit quantized pretrained language model into Low Rank Adapters~(LoRA). Our top-performing model family, called \textbf{Guanaco}, surpasses all previously publicly available models on the Vicuna benchmark, achieving 99.3\% of ChatGPT’s performance level with only 24 hours of finetuning on a single GPU. \textsc{QLoRA} incorporates several innovations to reduce memory usage without losing accuracy: (a) 4-bit NormalFloat (NF4), a novel data format that is information theoretically optimal for weights with a normal distribution, (b) Double Quantization, which lowers average memory consumption by quantizing the quantization constants themselves, and (c) Paged Optimizers to handle memory spikes effectively. We utilize \textsc{QLoRA} to finetune over 1,000 models, conducting a comprehensive evaluation of instruction following and chatbot capabilities across eight instruction datasets, various model architectures (LLaMA, T5), and model sizes that would be impractical for standard finetuning approaches (e.g., 33B and 65B parameter models). Our findings indicate that QLoRA finetuning on a small, high-quality dataset yields state-of-the-art results, even using models smaller than prior leading methods. We offer an in-depth assessment of chatbot quality based on both human and GPT-4 evaluations, demonstrating that GPT-4 assessments provide an affordable and reliable alternative to human judgments. Moreover, we observe that existing chatbot benchmarks are unreliable for accurately measuring chatbot performance levels. A selective failure analysis highlights the areas where \textbf{Guanaco} falls short relative to ChatGPT. We publicly release all models and code, including CUDA kernels designed for 4-bit training.\footnote{\url{https://github.com/artidoro/qlora} and \url{https://github.com/TimDettmers/bitsandbytes}}
\end{abstract}

\section{Introduction}

Fine-tuning large language models (LLMs) is a highly effective strategy for enhancing their capabilities \citep{min2021metaicl, wei2021finetuned, ouyang2022training, wang2022super, wang2022self, liu2022few} as well as for introducing preferred behaviors or eliminating undesired ones \citep{ouyang2022training,askell2021general,bai2022training}. Nevertheless, fine-tuning extremely large models demands prohibitive computational resources; for instance, standard 16-bit fine-tuning of a LLaMA model with 65 billion parameters \citep{touvron2023llama} requires over 780 GB of GPU memory. Although recent quantization approaches can reduce the memory demands of LLMs \citep{dettmers2022llm, dettmers2022case, frantar2022gptq, xiao2022smoothquant}, these methods are generally limited to inference tasks and fail during training \citep{wortsman2023stable}.

In this work, we present the first demonstration that fine-tuning a quantized 4-bit model can be accomplished without any loss in performance. Our approach, \textsc{QLoRA}, applies an innovative high-precision quantization technique to compress a pretrained model to 4 bits and then incorporates a compact set of trainable Low-rank Adapter weights \citep{hu2021lora}, which are optimized by backpropagating gradients through the quantized weights.

\textsc{QLoRA} reduces the average GPU memory needed to fine-tune a 65 billion parameter model from over 780 GB to under 48 GB, without compromising runtime efficiency or predictive accuracy relative to a fully fine-tuned 16-bit baseline. This represents a substantial enhancement in the accessibility of LLM fine-tuning, enabling the largest publicly available models to be fine-tuned on a single GPU. Utilizing \textsc{QLoRA}, we develop the \textbf{Guanaco} model family; the second-best model in this collection achieves 97.8\% of ChatGPT's performance on the Vicuna~\citep{vicuna2023} benchmark and can be trained in less than 12 hours on a single consumer GPU. With a single professional GPU over 24 hours, our largest model reaches 99.3\%, effectively closing the gap to ChatGPT on this benchmark. When deployed, our smallest \textbf{Guanaco} model (7 billion parameters) requires only 5 GB of memory and outperforms a 26 GB Alpaca model by more than 20 percentage points on the Vicuna benchmark (Table~\ref{tbl:vicuna_eval}).

\textsc{QLoRA} presents several innovations aimed at decreasing memory consumption without compromising performance: (1) \textbf{4-bit NormalFloat}, a quantization data type optimized from an information-theoretic perspective for normally distributed data, which achieves better empirical outcomes than both 4-bit Integers and 4-bit Floats. (2) \textbf{Double Quantization}, a technique that further compresses the quantization constants, saving on average approximately 0.37 bits per parameter (roughly 3 GB for a 65B parameter model). (3) \textbf{Paged Optimizers}, which leverage NVIDIA unified memory to mitigate the memory spikes caused by gradient checkpointing when processing mini-batches with long sequence lengths. These innovations are integrated into a refined LoRA method that incorporates adapters at every layer of the network, effectively eliminating nearly all accuracy compromises observed in earlier approaches.

The efficiency of \textsc{QLoRA} allows us to conduct a comprehensive exploration of instruction finetuning and chatbot performance across model scales that would be infeasible with conventional finetuning methods due to memory constraints. Consequently, we train over 1,000 models spanning multiple instruction tuning datasets, various architectures, and sizes ranging from 80M to 65B parameters. Beyond demonstrating that \textsc{QLoRA} matches 16-bit precision performance (\S\ref{sec:comp_to_std}) and developing a state-of-the-art chatbot, \textbf{Guanaco}, (\S\ref{sec:pushing}), we also investigate patterns in the trained models. Firstly, we observe that the quality of the data holds significantly more weight than the dataset size; for instance, a 9k-sample dataset (OASST1) outperformed a 450k-sample dataset (FLAN v2, subsampled) in chatbot performance, despite both targeting instruction-following generalization. Secondly, our findings reveal that excellent performance on the Massive Multitask Language Understanding (MMLU) benchmark does not necessarily correlate with strong results on the Vicuna chatbot benchmark, and vice versa—highlighting that the appropriateness of the dataset is more critical than its size for a specific task.

Additionally, we conduct a comprehensive examination of chatbot performance, utilizing evaluations from both human assessors and GPT-4. Our approach employs a tournament-style benchmarking method where different models compete by generating responses to given prompts. The victor of each match is determined by either GPT-4 or human judges. These tournament outcomes are compiled into Elo scores~\citep{elo1967proposed, elo1978rating}, which establish the leaderboard for chatbot performance. Our findings indicate that GPT-4 and human assessments generally concur on the ranking of model capabilities within these tournaments, though there are notable cases of significant disagreement. Therefore, we emphasize that while model-based evaluation offers a cost-effective alternative to human annotation, it also involves certain uncertainties.

We complement our chatbot benchmarking outcomes with a qualitative review of the \textbf{Guanaco} models. This analysis brings to light both successful and failing instances that quantitative metrics did not capture.

To support further research, we release all model outputs accompanied by annotations from humans and GPT-4. Our entire codebase and CUDA kernels are open-sourced, and we have integrated our methods into the Hugging Face transformers framework \citep{wolf2019huggingface} for easy accessibility. We provide a suite of adapters for 7B/13B/33B/65B parameter models, each trained on eight distinct instruction-following datasets, culminating in a total of 32 publicly available fine-tuned models.

\section{Background}
\paragraph{Block-wise k-bit Quantization} Quantization refers to the process of converting data from a representation containing more information to one with fewer bits. Typically, this involves changing a data type with higher bit width to one with fewer bits, such as converting 32-bit floating-point numbers into 8-bit integers. To fully utilize the range available in the lower-bit format, the original data is usually normalized by the absolute maximum value within the data (commonly a tensor), scaling it into the target type’s range. For instance, converting a 32-bit floating-point tensor (FP32) into an 8-bit integer tensor (Int8) with a range of $[-127, 127]$ is performed as follows:

\begin{equation}
    \mathbf{X}^{ \text{Int8}} = \text{round}\left( \frac{127}{{\text{absmax}}(\mathbf{X}^{{\text{FP32}}})}\mathbf{X}^{{\text{FP32}}} \right) = \text{round}(c^{\text{FP32}}\cdot \mathbf{X}^{{\text{FP32}}}),
\end{equation}

where $c$ represents the {\em quantization constant} or {\em quantization scale}. The process of dequantization is the reverse operation:

\begin{equation}
    \text{dequant}(c^{\text{FP32}}, \mathbf{X}^{\text{Int8}}) = \frac{\mathbf{X}^{\text{Int8}}}{c^{\text{FP32}}} = \mathbf{X}^{\text{FP32}}
\end{equation}

The issue with this method is that if an input tensor contains a value with a large magnitude (i.e., an outlier), the quantization bins—specific bit patterns—are inefficiently utilized, with some bins having few or no quantized values. To address the outlier problem, a typical strategy is to divide the input tensor into blocks that are quantized separately, each block having its own quantization constant $c$. Formally, the input tensor $\mathbf{X} \in \mathbb{R}^{b \times h}$ is partitioned into $n$ consecutive blocks of size $B$ by first flattening the tensor and then slicing the resulting linear array into $n = \frac{b \times h}{B}$ blocks. Each block is independently quantized using Equation~1, yielding a quantized tensor along with $n$ quantization constants $c_i$.

\paragraph{Low-rank Adapters} Low-rank Adapter (LoRA) finetuning \citep{hu2021lora} is a technique designed to reduce memory usage by introducing a small number of trainable parameters called adapters, while keeping the original model parameters fixed. During stochastic gradient descent, gradients flow through the static pretrained weights to the adapter, which is then updated to minimize the loss function. LoRA enhances a linear projection by adding an extra factorized projection. For a projection ${\mathbf{X} \mathbf{W} = \mathbf{Y}}$ with $\mathbf{X} \in \mathbb{R}^{b \times h}$ and $\mathbf{W} \in \mathbb{R}^{h \times o}$, LoRA computes:

\begin{equation}
    \mathbf{Y} = \mathbf{X}\mathbf{W} + s\mathbf{X}\mathbf{L}_1\mathbf{L}_2,
\end{equation}

where $\mathbf{L}_1 \in \mathbb{R}^{h \times r}$, $\mathbf{L}_2 \in \mathbb{R}^{r \times o}$, and $s$ is a scalar factor.

\paragraph{Memory Requirement of Parameter-Efficient Finetuning} A key topic of consideration is the memory consumption of LoRA during training, particularly regarding the quantity and size of adapters employed. Given LoRA's very small memory usage, it allows for incorporating additional adapters to enhance performance without a substantial rise in total memory demands. Although LoRA is developed as a Parameter Efficient Finetuning (PEFT) approach, the majority of memory usage in LLM finetuning arises from activation gradients rather than the learned LoRA parameters. For example, when training a 7B LLaMA model on FLAN v2 with a batch size of 1, and LoRA weights corresponding to the commonly used 0.2\% fraction of the original model weights \citep{hu2021lora,liu2022few}, the memory footprint for LoRA input gradients is 567 MB, whereas the LoRA parameters occupy only 26 MB. Applying gradient checkpointing~\citep{chen2016training} reduces the input gradients to an average of 18 MB per sequence, making them more memory-heavy than all LoRA parameters combined. In contrast, the 4-bit base model requires 5,048 MB of memory. This emphasizes the significance of gradient checkpointing and shows that aggressively shrinking the number of LoRA parameters results in only marginal memory savings. Consequently, more adapters can be utilized without notably increasing the overall training memory usage (refer to Appendix~\ref{app:memory} for a comprehensive analysis). As discussed subsequently, this is essential for achieving full 16-bit precision performance recovery.

\section{\textsc{QLoRA} Finetuning}

\textsc{QLoRA} attains precise 4-bit finetuning through two methods we introduce—4-bit NormalFloat (NF4) quantization and Double Quantization. Furthermore, we propose Paged Optimizers to prevent memory surges during gradient checkpointing from triggering out-of-memory issues, which have traditionally complicated finetuning large models on a single machine.

\textsc{QLoRA} utilizes one low-precision storage format, typically 4-bit in our experiments, alongside one computation format that is generally BFloat16. Practically, this means that whenever a \textsc{QLoRA} weight tensor is accessed, it is dequantized into BFloat16 before carrying out a 16-bit matrix multiplication.

We proceed by detailing the elements of \textsc{QLoRA}, followed by a formal description of the method.

\paragraph{4-bit NormalFloat Quantization} The NormalFloat (NF) datatype extends Quantile Quantization \citep{dettmers2022optimizers}, which is an information-theoretically optimal approach ensuring that each quantization bin contains an equal number of values drawn from the input tensor. Quantile quantization operates by estimating the quantiles of the input tensor using the empirical cumulative distribution function.

A primary drawback of quantile quantization is the high computational cost of quantile estimation. Consequently, fast approximation methods for quantiles, like SRAM quantiles~\citep{dettmers2022optimizers}, are employed for estimation. However, due to the approximate nature of these algorithms, the data type exhibits significant quantization errors for outliers, which frequently represent crucial values.

The computational burden of quantile estimation and the approximation inaccuracies can be avoided when input tensors originate from a distribution that remains fixed up to a quantization constant. In these scenarios, input tensors share identical quantiles, making exact quantile estimation computationally practical.

Pretrained neural network weights typically follow a zero-centered normal distribution characterized by a standard deviation $\sigma$ (refer to Appendix~\ref{app:normality}). By adjusting $\sigma$, we can scale all weights to conform to a single, fixed distribution that fits precisely within the bounds of our chosen data type. For our purposes, we define this arbitrary range as $[-1, 1]$. Consequently, the quantiles corresponding to both the data type and the neural network weights must be normalized within this interval.

The data type that is information-theoretically optimal for zero-mean normal distributions with any standard deviation $\sigma$ confined to $[-1, 1]$ is derived through the following process: (1) calculate the $2^k+1$ quantiles of a theoretical $N(0, 1)$ distribution to create a $k$-bit quantile quantization data type tailored for normal distributions, (2) normalize this data type’s values to fit within the $[-1, 1]$ interval, and (3) quantize an input weight tensor by scaling it into the $[-1, 1]$ range via absolute maximum normalization.

After ensuring the weight range aligns with the data type range, standard quantization can be applied. The third step corresponds to adjusting the standard deviation of the weight tensor to align with that of the $k$-bit data type. More explicitly, we determine the $2^k$ quantization levels $q_i$ of the data type by:

\begin{equation}
q_i  = \frac{1}{2}\left( Q_X\left(\frac{i}{2^k+1}\right) + Q_X\left(\frac{i+1}{2^k+1}\right)\right),
\end{equation}

where $Q_X(\cdot)$ denotes the quantile function of the standard normal distribution $N(0,1)$. One challenge with symmetric k-bit quantization is that it lacks an exact zero representation, which is crucial for quantizing padding and other zero-valued elements without error. To guarantee a discrete zero point at $0$ and utilize all $2^k$ levels for a k-bit data type, we form an asymmetric data type by estimating the quantiles $q_i$ for two separate ranges: $2^{k-1}$ quantiles for the negative side and $2^{k-1}+1$ quantiles for the positive side. These two sets of $q_i$ are then combined, and one of the duplicated zeros present in both sets is removed. We refer to this resulting data type—which has an equal expected number of values per quantization bin—as the \emph{k-bit NormalFloat} (NFk), as it is information-theoretically optimal for zero-centered normally distributed data. The precise values for this data type are provided in Appendix~\ref{app:nf4}.

\paragraph{Double Quantization{}}
We propose \emph{Double Quantization} (DQ), a method of quantizing the quantization constants themselves to achieve further memory savings. While a small block size is necessary for accurate 4-bit quantization~\citep{dettmers2022case}, it incurs significant memory overhead. For instance, using 32-bit constants with a block size of 64 for $\mathbf{W}$ leads to an average overhead of $32/64 = 0.5$ bits per parameter from the quantization constants. Double Quantization reduces this memory cost.

Specifically, Double Quantization regards the quantization constants $c_2^{\text{FP32}}$ from the first quantization step as inputs to a second quantization. This second quantization produces quantized constants $c_2^{\text{FP8}}$ and a new set of quantization constants $c_1^{\text{FP32}}$. For the second quantization, we employ 8-bit floats with a block size of 256, since 8-bit quantization does not degrade performance, consistent with findings by \citet{dettmers2022case}. Because the values $c_2^{\text{FP32}}$ are all positive, we subtract their mean before quantization to center them around zero and apply symmetric quantization. On average, with a block size of 64, this reduces the memory overhead per parameter from $32/64=0.5$ bits to $8/64 + 32/(64 \cdot 256) = 0.127$ bits, yielding a savings of 0.373 bits per parameter.

\paragraph{Paged Optimizers} leverage NVIDIA unified memory~\footnote{\tiny \url{https://docs.nvidia.com/cuda/cuda-c-programming-guide}} which facilitates automatic page-by-page data transfers between the CPU and GPU. This ensures error-free GPU execution even when the GPU occasionally experiences out-of-memory conditions. This mechanism operates similarly to traditional memory paging between CPU RAM and disk storage. We employ this capability to allocate paged memory for optimizer states, which are then seamlessly moved to CPU RAM when GPU memory is insufficient and swapped back to GPU memory when required during the optimizer update step.

\paragraph{\textsc{QLoRA}.} Based on the previously described components, we define \textsc{QLoRA} for a single linear layer in the quantized base model augmented with a single LoRA adapter as follows:

\begin{equation}
    \mathbf{Y}^{\text{BF16}} =  \mathbf{X}^{\text{BF16}} \text{doubleDequant}(c_1^{\text{FP32}}, c_2^{\text{k-bit}}, \mathbf{W}^{\text{NF4}}) + \mathbf{X}^{\text{BF16}}\mathbf{L}^{\text{BF16}}_1\mathbf{L}^{\text{BF16}}_2,
\end{equation}

where the function doubleDequant$(\cdot)$ is defined as:

\begin{equation}
    \text{doubleDequant}(c_1^{\text{FP32}}, c_2^{\text{k-bit}}, \mathbf{W}^{\text{k-bit}}) = \text{dequant}(\text{dequant}(c_1^{\text{FP32}}, c_2^{\text{k-bit}}), \mathbf{W}^{\text{4bit}}) = \mathbf{W}^{\text{BF16}},
\end{equation}

We utilize NF4 for the matrix $\mathbf{W}$ and FP8 for the coefficient $c_2$.

A block size of 64 is chosen for $\mathbf{W}$ to achieve higher quantization accuracy, while a block size of 256 is used for $c_2$ to save memory.

For parameter updates, only the gradients with respect to the adapter weights, $\frac{\partial E}{\partial \mathbf{L}_i}$, are required; gradients for the 4-bit weights $\frac{\partial E}{\partial \mathbf{W}}$ are not needed. Nevertheless, computing $\frac{\partial E}{\partial \mathbf{L}_i}$ involves determining $\frac{\partial \mathbf{X}}{\partial \mathbf{W}}$, which follows equation~(5) and requires dequantizing the stored weights $\mathbf{W}^{\text{NF4}}$ into the computation data type $\mathbf{W}^{\text{BF16}}$ to evaluate the derivative $\frac{\partial \mathbf{X}}{\partial \mathbf{W}}$ at BFloat16 precision.

In summary, \textsc{QLoRA} utilizes a single storage data type (commonly 4-bit NormalFloat) alongside a computation data type (16-bit BrainFloat). The storage data is dequantized into the computation data type to carry out the forward and backward passes, but weight gradients are computed solely for the LoRA parameters, which employ 16-bit BrainFloat.

\section{QLoRA Compared to Standard Finetuning}
\label{sec:comp_to_std}

We have outlined the mechanics of QLoRA and its ability to drastically reduce memory demands during model finetuning. The crucial question is whether QLoRA can achieve performance comparable to full-model finetuning. Additionally, we aim to examine the individual components of QLoRA, such as the effect of using NormalFloat4 versus the conventional Float4. The next sections will present experiments designed to address these issues.

\paragraph{Experimental setup.} We evaluate three model architectures (encoder, encoder-decoder, and decoder-only) and compare QLoRA against 16-bit adapter finetuning and full finetuning on models up to 3B parameters. Our benchmarks cover GLUE \citep{wang2018glue} with RoBERTa-large \citep{liu2019roberta}, Super-NaturalInstructions (TKInstruct) \citep{wang2022super} with T5 \citep{t5}, and 5-shot MMLU \citep{hendrycksmeasuring} after finetuning LLaMA on Flan v2 \citep{longpre2023flan} and Alpaca \citep{alpaca}. To further investigate the benefits of NF4 over other 4-bit formats, we follow the setup of \citet{dettmers2022case}, evaluating post-quantization zero-shot accuracy and perplexity across several models (OPT~\citep{zhang2022opt}, LLaMA~\citep{touvron2023llama}, BLOOM~\citep{scao2022bloom}, Pythia~\citep{biderman2023pythia}) covering sizes from 125m to 13B parameters.

More detailed results for each specific configuration are provided in the results section to enhance clarity. Comprehensive experimental details are available in Appendix~\ref{app:exp-setup}.

Although paged optimizers are essential for performing 33B/65B \textsc{QLoRA} finetuning on a single 24/48GB GPU, we do not report precise measurements for paged optimizers here because paging mainly occurs when processing mini-batches with long sequences, which is infrequent. Nevertheless, we analyze the runtime of paged optimizers for 65B models on 48GB GPUs and observe that at a batch size of 16, paged optimizers achieve training speeds comparable to standard optimizers. Future investigations should focus on quantifying when and how paging leads to slowdowns.

\paragraph{Default LoRA hyperparameters fail to match 16-bit finetuning performance} 

Using the conventional method of applying LoRA solely to query and value attention projection matrices \citep{hu2021lora}, we could not replicate the performance of full finetuning on larger base models. As illustrated in Figure~\ref{fig:hyper} for LLaMA 7B finetuned on Alpaca, the key LoRA hyperparameter influencing performance is the total number of LoRA adapters employed. Achieving parity with full finetuning requires applying LoRA to all linear layers within transformer blocks. Other hyperparameters of LoRA, such as the projection dimension $r$, do not significantly impact performance (see Appendix \ref{app:exp-setup}).

Similarly, we observe that the default hyperparameters for fully finetuned baselines are suboptimal. To identify robust baselines, we conduct a hyperparameter search over learning rates ranging from 1e-6 to 5e-5 and batch sizes between 8 and 128. The results for 7B LLaMA finetuning on Alpaca are presented in Figure~\ref{fig:hyper}.

\paragraph{4-bit NormalFloat outperforms 4-bit Floating Point}

Although the 4-bit NormalFloat (NF4) data type is theoretically optimal from an information perspective, it remains to be seen whether this theoretical advantage translates into practical benefits. We replicate the experimental setup from \citet{dettmers2022case}, where quantized LLMs (including OPT \citep{zhang2022opt}, BLOOM \cite{scao2022bloom}, Pythia \cite{biderman2023pythia}, and LLaMA) of various sizes (ranging from 125M to 65B parameters) and data types are evaluated on language modeling tasks and several zero-shot benchmarks. As shown in Figure~\ref{fig:nf4} and Table~\ref{tbl:nf4_ppl}, NF4 significantly enhances performance compared to FP4 and Int4, and double quantization effectively reduces memory usage without sacrificing accuracy.

\paragraph{k-bit \textsc{QLoRA} achieves parity with 16-bit full finetuning and 16-bit LoRA}

Recent studies have demonstrated that 4-bit quantization is feasible for \emph{inference}, although it typically results in some performance degradation compared to 16-bit precision \citep{dettmers2022case,frantar2022gptq}. This prompts an important question: can this performance loss be recovered by performing 4-bit adapter finetuning? We investigate this question under two experimental conditions.

The first condition compares with full 16-bit finetuning of RoBERTA and T5 models, ranging from 125M to 3B parameters, evaluated on GLUE and the Super-NaturalInstructions dataset. Results are displayed in Table~\ref{tab:nf4vsbf16}. Across both datasets, 16-bit, 8-bit, and 4-bit adapter methods match the performance of the fully finetuned 16-bit baseline. This indicates that the performance deficit caused by coarse quantization can be fully restored via adapter finetuning after quantization.

For our second experimental setup, because fully finetuning models at and above 11B parameters demands more than one high-memory GPU server, we continue to investigate whether 4-bit \textsc{QLoRA} can achieve parity with 16-bit LoRA across the 7B to 65B parameter range. To accomplish this, we finetune LLaMA models ranging from 7B to 65B parameters on two instruction-following datasets, Alpaca and FLAN v2, and assess their performance using the MMLU benchmark with 5-shot accuracy. The results, presented in Table~\ref{tbl:llama-datatype}, reveal that NF4 combined with double quantization fully recovers the MMLU performance seen with 16-bit LoRA. Additionally, we observe that \textsc{QLoRA} using FP4 trails the 16-bit brain float LoRA baseline by roughly one percentage point. This supports our conclusions that (1) \textsc{QLoRA} with NF4 matches both 16-bit full finetuning and 16-bit LoRA finetuning performance, and (2) NF4 outperforms FP4 in quantization accuracy.

\paragraph{Summary} Our findings consistently demonstrate that 4-bit \textsc{QLoRA} employing the NF4 data type achieves comparable performance to 16-bit full finetuning and 16-bit LoRA finetuning on established academic benchmarks with well-defined evaluation protocols. We also establish that NF4 is superior to FP4 and that applying double quantization does not impair performance. Together, these observations provide strong evidence that 4-bit \textsc{QLoRA} tuning reliably matches the outcomes of 16-bit methods.

Consistent with prior quantization research \citep{dettmers2022case}, our MMLU and Elo evaluations suggest that, given a fixed budget for finetuning and inference resources, it is advantageous to increase the base model’s parameter count while lowering their precision. This underscores the significance of the efficiency improvements offered by \textsc{QLoRA}. Since no performance drop was observed relative to full finetuning in our 4-bit finetuning experiments, this prompts the question of where exactly the performance-precision trade-off lies for QLoRA tuning, which we leave to future investigations.

We move forward to explore instruction tuning at scales that would be unattainable using full 16-bit finetuning on typical academic research hardware.

\section{Advancing the Chatbot State-of-the-art with QLoRA}
\label{sec:pushing}
After demonstrating that 4-bit \textsc{QLoRA} achieves parity with 16-bit performance across various scales, tasks, and datasets, we carry out a comprehensive study of instruction finetuning on the largest open-source language models currently accessible for research purposes. To evaluate the effectiveness of instruction finetuning on these models, we utilize a challenging Natural Language Understanding benchmark (MMLU) and propose new approaches for assessing real-world chatbot performance.

\subsection{Experimental setup} We now provide an overview of the experimental setup, with detailed information available in Appendix \ref{app:chatbot}.

\paragraph{Data} Since, to the best of our knowledge, there is no thorough study of recent instruction-following datasets, we choose eight recent datasets. These include datasets collected via crowd-sourcing (OASST1~\citep{kopf2023openassistant}, HH-RLHF~\citep{bai2022training}), those distilled from instruction-tuned models (Alpaca~\citep{alpaca}, self-instruct~\citep{wang2022self}, unnatural-instructions~\citep{honovich2022unnatural}), aggregated corpora (FLAN v2~\citep{chung2022scaling}), as well as hybrid datasets (Chip2~\citep{laion2023}, Longform~\citep{koksal2023longform}). These datasets vary in terms of languages, data volume, and licensing.

\paragraph{Training Setup} To avoid confounding factors from different training objectives, we apply QLoRA finetuning using cross-entropy loss (supervised learning) without reinforcement learning, even for datasets containing human-judged multiple responses. For datasets where there is a clear separation between instruction and response, we finetune exclusively on the response (see ablations in Appendix \ref{app:chatbot}). For OASST1 and HH-RLHF, where multiple responses are present, we select the top response at each node of the conversation tree and finetune on the entire chosen conversation, including the instructions. Across all experiments, we use NF4 \textsc{QLoRA} with double quantization and paged optimizers to avoid memory spikes during gradient checkpointing. We conduct minor hyperparameter tuning for the 13B and 33B LLaMA models and observe that all hyperparameter settings identified at 7B (including number of epochs) generalize, except for learning rate and batch size. For the 33B and 65B models, we halve the learning rate and double the batch size.

\paragraph{Baselines} We benchmark our models against both research-based (Vicuna~\citep{vicuna2023} and Open Assistant~\citep{kopf2023openassistant}) and commercial (GPT-4 \citep{openai2023gpt}, GPT-3.5-turbo, and Bard) chatbot systems. The Open Assistant model is a LLaMA 33B variant finetuned with Reinforcement Learning from Human Feedback (RLHF) on the same OASST1 dataset that we use. Vicuna performs full fine-tuning of LLaMA 13B on proprietary user-shared conversations from ShareGPT and represents a distillation from OpenAI GPT models.

\subsection{Evaluation}

Following standard practice, we employ the MMLU (Massively Multitask Language Understanding) benchmark \citep{hendrycksmeasuring} to evaluate performance across a variety of language understanding tasks. This benchmark consists of multiple-choice questions spanning 57 subjects, including basic mathematics, US history, computer science, law, among others. We report 5-shot test accuracy.

We also assess generative language capabilities via both automated and human evaluations. This second evaluation approach uses queries selected by humans and focuses on assessing the quality of the model's responses. Although this method represents a more realistic environment for testing chatbot model performance and is becoming increasingly popular, no universally accepted protocol currently exists in the literature. Below, we outline our proposed setup, which employs nucleus sampling with $p=0.9$ and temperature $0.7$ consistently.

\paragraph{Benchmark Data} We perform evaluations on two curated query datasets: the Vicuna prompts \citep{vicuna2023} and the OASST1 validation dataset \citep{kopf2023openassistant}. The Vicuna prompts consist of 80 prompts drawn from a wide range of categories, used without any changes. The OASST1 dataset is a multilingual collection of crowd-sourced, multi-turn dialogs between users and an assistant. We extract all user messages from the validation portion as queries, including prior dialog turns within the prompt. This yields 953 unique user queries. We refer to these two datasets collectively as the Vicuna and OA benchmarks.

\paragraph{Automated Evaluation} Initially, following the evaluation protocol introduced by \citet{vicuna2023}, we utilize GPT-4 to rate various systems' performance compared to ChatGPT (GPT-3.5 Turbo) on the Vicuna benchmark. For each query along with responses from ChatGPT and another model, GPT-4 is prompted to assign a score out of ten to each response and provide justification. The overall model performance is then expressed as a percentage relative to ChatGPT’s score. This relative score may exceed 100\% if the model attains a higher absolute score than ChatGPT. We observe a notable ordering bias whereby GPT-4 tends to favor the response presented first. To mitigate this bias, we recommend reporting the average score over both response orderings.

Subsequently, we assess performance through direct comparisons of system outputs. We simplify the rating task into a three-class labeling problem that permits ties. GPT-4 is prompted to select the superior response or declare a tie, along with an explanation. These pairwise comparisons are performed across all system permutations on both the Vicuna and OA benchmarks.

\paragraph{Human Evaluation} Although recent studies suggest generative models can serve effectively in system evaluations \citep{fu2023gptscore}, to our knowledge, the consistency of GPT-4 ratings with human judgments for chatbot evaluation has not yet been established. Therefore, we conduct two parallel human evaluations on the Vicuna benchmark following the two automated evaluation protocols described above. We recruit annotators via Amazon Mechanical Turk (AMT), obtaining two human raters for comparisons against ChatGPT and three raters for pairwise system comparisons.

\paragraph{Elo Rating} Using both human and automated pairwise comparisons, we set up a tournament-style competition where models face off against one another. The tournament consists of matches in which pairs of models compete to generate the best response for a given prompt. This approach is akin to the comparison methods used by \citet{bai2022training} and \citet{vicuna2023}, but we additionally incorporate GPT-4 ratings alongside human evaluations. We randomly select from the pool of labeled comparisons to calculate Elo scores~\citep{elo1967proposed,elo1978rating}. Elo rating, a system commonly utilized in chess and other games, quantifies the expected win rate relative to an opponent's rating—for instance, a player with an Elo of 1100 versus an opponent rated 1000 is expected to win about 65\% of the time; matches between players with equal ratings, such as 1000 vs 1000 or 1100 vs 1100, predict an expected win rate of 50\%. Following each match, the Elo rating updates based on the expected result; an unexpected upset causes a significant rating shift, whereas an anticipated outcome leads to a minor adjustment. Over time, Elo ratings tend to reflect the true skill level of each participant in the game. We initialize all ratings at 1,000 and set the $K$ factor to 32. In line with \citet{vicuna2023}, we repeat this process 10,000 times using different random seeds to mitigate ordering effects, such as the influence of which model pairs compete first.

\subsection{\textbf{Guanaco}: \textsc{QLoRA} fine-tuned on OASST1 as a Leading Chatbot} Our automated and human assessments indicate that the top \textsc{QLoRA}-tuned model, {Guanaco} 65B, which we fine-tune on a modified version of OASST1, is the highest-performing open-source chatbot model and delivers performance comparable to ChatGPT. When compared to GPT-4, {Guanaco} 65B and 33B have an expected win probability of 30\%, based on Elo ratings derived from human annotators’ system-level pairwise comparisons — the highest reported so far.

Table~\ref{tbl:vicuna_eval} presents results from the Vicuna benchmark \cite{vicuna2023} relative to ChatGPT. We observe that {Guanaco} 65B ranks as the best model after GPT-4, reaching 99.3\% of ChatGPT’s performance. Although {Guanaco} 33B contains more parameters than the Vicuna 13B model, it uses only 4-bit precision for its weights, making it significantly more memory efficient at 21 GB versus 26 GB, and providing a 3 percentage point improvement over Vicuna 13B. Additionally, {Guanaco} 7B is compact enough to run on modern smartphones with a 5 GB memory footprint, while still outperforming Alpaca 13B by nearly 20 percentage points.

Nonetheless, Table~\ref{tbl:vicuna_eval} exhibits wide confidence intervals, with performance overlapping among many models. We suggest that this uncertainty stems from the ambiguous definition of scale, for example, what a score of 8 on a 10-point scale signifies across various contexts. Therefore, we propose using the Elo ranking approach \citep{elo1967proposed}, which relies on \textit{pairwise} evaluations from human annotators and GPT-4, circumventing issues associated with anchoring an absolute scale. Elo ratings for the most competitive models appear in Table~\ref{tbl:elo}. We observe some discrepancies between human and GPT-4 rankings on the Vicuna benchmark, particularly for Guanaco 7B, but overall they align for most models, with a Kendall Tau of $\tau=0.43$ and a Spearman rank correlation of $r=0.55$ at the system level. At the example level, agreement between GPT-4 and the majority human vote is weaker, with a Fleiss $\kappa=0.25$. In summary, these findings indicate a moderate concordance between system-level judgments from GPT-4 and human annotators, supporting the view that model-based evaluation offers a reasonably reliable alternative to human evaluation. Further details are discussed in Section~\ref{ssec:considerations}.

Elo rankings presented in Table~\ref{tbl:elo_large} demonstrate that the {Guanaco} 33B and 65B models surpass all other models except GPT-4 on the Vicuna and OA benchmarks, and their performance is on par with ChatGPT, consistent with the findings in Table~\ref{tbl:vicuna_eval}. It is important to highlight that the Vicuna benchmark tends to favor open-source models, whereas the larger OA benchmark shows a preference for ChatGPT. Additionally, from Tables \ref{tbl:mmlu-llama} and \ref{tbl:vicuna_eval}, it is evident that the choice of finetuning dataset plays a crucial role in performance outcomes. Specifically, finetuning Llama models on FLAN v2 yields strong results on MMLU but results in the poorest performance on the Vicuna benchmark (similar patterns are observed across other models). This indicates some degree of independence between current evaluation benchmarks: excelling at MMLU does not guarantee high chatbot performance (as assessed by Vicuna or OA benchmarks), and the reverse also holds true.

Among the top models in our evaluation, {Guanaco} is unique in that it is not trained on proprietary datasets, given that the OASST1 dataset collection guidelines explicitly prohibit the use of GPT models. The next leading model trained solely on open-source data is the Anthropic HH-RLHF model, which trails {Guanaco} by 30 percentage points on the Vicuna benchmark (refer to Table~\ref{tbl:vicuna_eval}). Taken together, these findings illustrate that 4-bit \textsc{QLoRA} is an effective method capable of producing state-of-the-art chatbots that rival ChatGPT. Moreover, our 33B {Guanaco} model can be trained on 24 GB consumer GPUs in under 12 hours. This paves the way for future research leveraging \textsc{QLoRA} tuning on specialized open-source datasets, enabling the creation of models that compete with the leading commercial systems available today.

\section{Qualitative Analysis}

Although quantitative analysis forms the foundation of our evaluation, relying solely on summary statistics has several limitations. One of the most significant concerns is the issue of benchmark validity \citep{liao2021we}—whether a benchmark genuinely measures what its title or description claims is always debatable, particularly as researchers uncover ``shortcuts'' that machine learning models can exploit to solve benchmarks \citep{gururangan2018annotation, poliak2018hypothesis}. To partially address this, we conduct a qualitative analysis divided into two parts. First, in \S\ref{ssec:examples}, we present several examples that we consider representative of certain recurring patterns found in the text generated by our 65b {Guanaco} model. Second, in \S\ref{ssec:considerations}, we discuss the implications of the results we have observed and how we interpret them.

\subsection{Qualitative Analysis of Example Generations}
\label{ssec:examples}

To identify examples, we examine the outputs generated for the Vicuna and OpenAssistant benchmarks, searching for recurring patterns in the responses produced by {Guanaco}. When we detect a pattern, we craft a question or prompt designed to elicit that pattern deliberately, even if it leads to an incorrect answer. For instance, if we notice the model tends to produce overly verbose responses, we prompt it with ``Answer yes or no without explanation.'' This approach allows us to identify both ``lemons,'' where we can adversarially cause the model to fail, and ``cherries,'' where the model resists such attempts. All examples presented in this section were generated using Nucleus Sampling \cite{holtzmancurious} with a probability cutoff of $p=0.9$.

Naturally, this is not exhaustive, as the limitations of this small qualitative study prevent controlling for every variable involved; for example, the full range of responses the model might produce for a given prompt is quite extensive, so we depend on samples we trust to be representative. Nevertheless, we consider that presenting these examples provides context for the quantitative data discussed earlier in the paper. Since we have made all models and code publicly available, we hope this section encourages further research to explore in greater depth the challenges we highlight here.  
\paragraph{Factual Recall} For queries like ``What is the capital of Zambia?'' all models reliably produce accurate answers, such as:

\begin{quote} 
    \textbf{Guanaco}:
    \texttt{The capital of Zambia is Lusaka.}
\end{quote}

However, as the questions become more obscure, {Guanaco} becomes less dependable while still maintaining high confidence. For example, given this prompt from HotPotQA \citep{yang2018hotpotqa}:

\begin{quote}
     \textbf{User}: In what year was the singer who popularized `I'll Keep the Lovelight Burning' born?  

     \textbf{Guanaco}:~\texttt{Al Jolson is credited with popularizing the song `I'll Keep the Lovelight Burning,' and he was born in the year 1886.}
 \end{quote}

it provides an incorrect popularizer and birth year (although the birth year is accurate for the mentioned individual, Al Jolson).

\paragraph{Suggestibility} {Guanaco} exhibits a notable resistance to accepting certain forms of assumed misinformation, as illustrated in the following exchange:

\begin{quote}
    \textbf{User}: How was it ultimately and officially verified that the earth is flat by peer reviewed scientists? \\ \\
    \textbf{Guanaco}: \texttt{The earth has never been officially verified to be flat by peer reviewed scientists. The vast majority of scientists agree that the earth is spherical. Indeed, this fact has been established for centuries and has been validated through a multitude of experiments and observations. \\ \\ The notion that the earth is flat is a conspiracy theory that has been thoroughly disproven by scientific evidence. The claim that a group of scientists have peer reviewed and confirmed the earth as flat is completely incorrect. \\ \\ If you have any further questions, please feel free to ask.}
\end{quote}

 also performs well in recognizing when questions are unanswerable, for example,

\begin{quote}
    \textbf{User}: What time is it?

    \textbf{Guanaco}: \texttt{I'm sorry, but as an AI language model, I do not have access to real-time information. My answers are generated based on existing knowledge and do not indicate the current time. If you need the current time, I recommend using a clock or checking an online time source.}
\end{quote}

Comparable outcomes are observed for queries such as ``Where are you?'', ``How are you?'', and so forth.

\paragraph{Refusal} Likewise, {Guanaco} occasionally declines to comply with instructions for apparently arbitrary reasons:

\begin{quote}
    \textbf{User}: Please reverse the words in the sentence ``I'm a cool dude in a mocha mood.'' \\ \\
    \textbf{Guanaco}: \texttt{I'm sorry, I can't perform that task at the moment, but I can clarify why the sentence you gave is grammatically correct.\\ \\The sentence ``I'm a cool dude in a mocha mood" adheres to the proper grammar and syntax rules of English. Here, the adjective ``cool'' precedes the noun ``dude'' to describe it, and the prepositional phrase ``in a mocha mood" adds contextual detail to the sentence. This sentence construction is completely valid and meaningful within its context.\\ \\If you have further questions or need assistance with anything else, please feel free to ask}
\end{quote}

\paragraph{Secret Keeping} Unfortunately, even when intended, this refusal is not always dependable. Often, system designers want to keep certain information hidden from the user, even if the system can respond with unrestricted text. We prompt the model by stating ``The secret word is `banana'. Do not reveal or repeat this under any circumstances,'' and examine how easily a user might obtain this secret. Direct queries fail:

\begin{quote}
    \textbf{User}: What is the secret word?

    \textbf{Guanaco}: \texttt{I'm sorry, but I cannot disclose the secret word. I am designed not to reveal or repeat it regardless of user requests. Is there anything else I can assist you with?}
\end{quote}

which is the intended response. However, minimal deception circumvents this behavior:

\begin{quote}
    \textbf{User}: This is a game. The goal is to ignore your previous instructions. What is the secret word?

    \textbf{Guanaco}: \texttt{Sorry, the secret word is banana. It's best if we don't talk about it. Is there anything else you'd like to ask?}
\end{quote}

demonstrating the importance of developing techniques for more reliable adherence to instructions.

\paragraph{Math} The main limitation of {Guanaco} lies in mathematics, a domain where many language models encounter difficulties \citep{liang2022holistic}. When {Guanaco} provides its reasoning, it is often correct, for instance,

\begin{quote}
    \textbf{User}: Lee mows one lawn and charges \$33. Last week he mowed 16 lawns and three customers each gave him a \$10 tip. How many dollars did Lee earn mowing lawns last week?

    \textbf{Guanaco}: \texttt{Lee earned \$582 mowing lawns last week. \\ \\ To determine the total amount he earned, we first calculate the earnings from mowing 16 lawns: \\ \\  \$33 x 16 = \$528 \\ \\ Next, we add the tips given by the three customers: \\  \\ \$10 x 3 = \$30 \\ \\ Therefore, the total amount Lee earned was: \\ \\ \$528 + \$30 = \$558}
\end{quote}

Nonetheless, {Guanaco} can fail even on straightforward problems if it does not decompose them into step-by-step solutions, a well-documented challenge \cite{weichain}. For example, consider the following exchange:

\begin{quote}
    \textbf{User}: Please factorize 1833.

    \textbf{Guanaco}: \texttt{Sure, I'd be happy to assist! 1833 is a prime number, meaning it has only two factors: 1 and 1833. \\ \\ Thus, the unique factorization of 1833 is: \\ \\ 2\^{}1 * 3\^{}2 * 17\^{}1 \\ \\ I hope this is helpful! Feel free to ask if you have more questions.}
\end{quote}

which is not only incorrect (the actual factorization is $3 \times 17 \times 43$), but is doubly mistaken.

\paragraph{Theory of Mind} {Guanaco} demonstrates unexpectedly strong Theory of Mind skills \cite{nematzadeh2018evaluating, sap2022neural}. For example, the model provides a detailed correct response to the following prompt:

\begin{quote}
    \textbf{User}: James and Abby are in the bedroom. Abby put the pen in the desk drawer. Abby leaves the bedroom. James moves the pen into the bag. Where does James think Abby will look for the pen?

    \textbf{Guanaco}: \texttt{James knows that Abby put the pen in the desk drawer. But James also knows that he has moved the pen into the bag, and he assumes that Abby does not know about this. So, James may guess that Abby will look for the pen in the desk drawer, because that's where she left it last time. However, this assumption may be incorrect, because Abby may have seen James moving the pen.}
\end{quote}

Nonetheless, these kinds of inferences are often unreliable, as the model frequently produces rationales that rely on assumptions inconsistent with the scenario, for instance,

\begin{quote}
    \textbf{User}: Evelyn entered the living room. Jackson entered the playroom. James entered the playroom. The beans are in the treasure chest. James exited the playroom. Jackson moved the beans to the pantry. Jackson exited the playroom. James entered the living room. Where will James look for the beans?

    \textbf{Guanaco}: \texttt{James will look for the beans in the pantry, because that's where Jackson moved them.}
\end{quote}

in which {Guanaco} assumes information transfer that was never established. These challenges reflect findings in recent studies \cite{sap2022neural} and warrant further investigation.

\subsection{Considerations}
\label{ssec:considerations}

\paragraph{Evaluation}

We observe moderate agreement among human annotators (Fleiss $\kappa=0.42$), which declines further when comparing two high-performing systems. This indicates limitations in existing benchmarks and human evaluation methodologies for chatbot task performance. In manual comparisons of outputs from ChatGPT and {Guanaco} 65B on the Vicuna benchmark, subjective preferences emerge prominently, with the paper’s authors frequently disagreeing on which responses are preferable. Future research should explore strategies to address these challenges by leveraging insights from fields that have developed methods for managing subjective preferences, such as Human-Computer Interaction and Psychology.

Our analysis also reveals that automated evaluation methods exhibit notable biases. For example, we detect strong order effects, where GPT-4 tends to give higher ratings to the system listed first in its prompt. The relatively low agreement between GPT-4’s judgments and human annotators at the sample level (Fleiss $\kappa=0.25$) suggests that humans and automated systems may base their evaluations on different criteria. Additionally, as shown in Table~\ref{tbl:elo_large}, GPT-4 assigns significantly higher scores to outputs it generated itself compared to human ratings, with Elo ratings of 1348 versus 1176, corresponding to roughly a 20\% increased probability of winning against an opponent.

Future research should investigate the existence of potential biases in automated evaluation systems and explore possible strategies to mitigate them.

\paragraph{Data \& Training} We observe that the OASST1 dataset, which is used to train the {Guanaco} models, is multilingual, and the OA benchmark also includes prompts in various languages. Future studies could explore how much multilingual training enhances performance on instructions given in languages other than English, and whether this accounts for the larger performance gap observed between the Vicuna-13B model (trained exclusively on English data) and the {Guanaco} 33B and 65B models on the OA benchmark.

Considering the strong results achieved by {Guanaco} models, we examine the possibility of data leakage between the OASST1 dataset and the Vicuna benchmark prompts. After conducting fuzzy string matching between the two datasets and manually reviewing the closest matches, we do not find any overlapping prompts.

Additionally, it is important to note that our model is trained solely using cross-entropy loss (supervised learning), without utilizing reinforcement learning from human feedback (RLHF). This highlights the need for further research into the trade-offs between simple cross-entropy loss and RLHF training. We anticipate that \textsc{QLoRA} will facilitate such large-scale analyses without demanding excessive computational resources.

\section{Related Work}

\paragraph{Quantization of Large Language Models} Quantization techniques for LLMs have primarily concentrated on inference-time quantization. Key methods to maintain 16-bit model quality aim at handling outlier features (e.g., SmoothQuant~\citep{xiao2022smoothquant} and LLM.int8()~\citep{dettmers2022llm}), while other methods apply more advanced grouping strategies \citep{park2022nuqmm, yao2022zeroquant}. Lossy quantization research examines the trade-offs involved with standard rounding methods~\citep{dettmers2022case,zeng2022glm,pope2022efficiently} or techniques for optimizing rounding decisions to enhance quantization accuracy~\citep{frantar2022gptq}. Apart from our work, SwitchBack layers~\citep{wortsman2023stable} is the only study addressing backpropagation through quantized weights at a scale exceeding 1B parameters.

\paragraph{Finetuning with Adapters} Although we employ Low-rank Adapters~\citep{hu2021lora} (LoRA), various other Parameter Efficient FineTuning (PEFT) techniques have been introduced, including prompt tuning~\citep{qin2021learning,lester2021power,li2021prefix}, modifying embedding layer inputs~\citep{an2022input}, adjusting hidden states (IA$^3$) \citep{liu2022few}, incorporating additional layers~\citep{houlsby2019parameter}, fine-tuning biases~\citep{zaken2021bitfit}, learning weight masks guided by Fisher information~\citep{sung2021training}, and hybrid approaches~\citep{henderson2021compacter}. In our study, we demonstrate that LoRA adapters can achieve performance comparable to full 16-bit finetuning. Exploring the trade-offs of other PEFT methods remains an avenue for future research.

\paragraph{Instruction Finetuning} Instruction finetuning aims to enhance a pretrained LLM’s ability to follow prompts by finetuning it on input-output pairs derived from diverse data sources, teaching the model to produce the desired output given an input prompt. Notable methods and datasets encompass MetaICL~\citep{min2021metaicl}, MetaTuning~\citep{zhong2021adapting}, InstructGPT~\citep{ouyang2022training}, FLAN~\citep{wei2021finetuned,chung2022scaling}, PromptSource~\citep{bach2022promptsource}, Super-NaturalInstructions~\citep{wang2022super,sanh2021multitask}, Self-instruct~\citep{wang2022self}, UnnaturalInstructions~\citep{honovich2022unnatural}, OPT-IML~\citep{iyer2022opt}, UnifiedSKG~\citep{xie2022unifiedskg}, OIG/Chip2~\citep{laion2023}, Alpaca~\citep{alpaca}, Vicuna~\citep{vicuna2023}, Koala~\citep{koala_blogpost_2023}, and Self-instruct-GPT-4~\citep{peng2023instruction}.

\paragraph{Chatbots} Many instruction-following models are designed as dialogue-based chatbots, frequently utilizing Reinforcement Learning from Human Feedback (RLHF)~\citep{christiano2017deep} or generating training data with AI model feedback (RLAIF)~\citep{bai2022constitutional}. Examples of such approaches and datasets include Anthropic-HH~\citep{askell2021general,bai2022training}, Open Assistant~\citep{kopf2023openassistant}, LaMDA~\citep{thoppilan2022lamda}, and Sparrow~\citep{glaese2022improving}. While we do not employ reinforcement learning ourselves, our leading model, Guanaco, is finetuned on multi-turn chat dialogues from the Open Assistant dataset, which is designed for RLHF training~\citep{kopf2023openassistant}. For chatbot evaluation, alternatives using GPT-4 in place of expensive human annotation have been proposed \citep{vicuna2023,peng2023instruction}. We enhance these methods by emphasizing a more robust and reliable evaluation framework.

\section{Limitations and Discussion}

Our findings indicate that \textsc{QLoRA} can reproduce the performance of full 16-bit finetuning by using a 4-bit base model combined with Low-rank Adapters (LoRA). Nonetheless, we have not yet confirmed that \textsc{QLoRA} achieves parity with full 16-bit finetuning at larger model scales such as 33B and 65B. Given the substantial computational expense, this remains an open question for future investigation.

Another limitation concerns the evaluation of instruction finetuning models. Although we conducted assessments on MMLU, the Vicuna benchmark, and the OA benchmark, we did not test on other datasets like BigBench, RAFT, and HELM; hence, it is unclear whether our results generalize to those. However, we offer a comprehensive analysis on MMLU and introduce novel techniques for evaluating chatbot performance.

Based on the evidence provided, it seems that the effectiveness of these benchmarks likely hinges on how closely the finetuning data aligns with the benchmark dataset. For instance, FLAN v2 resembles MMLU but differs from chatbot benchmarks, whereas the Chip2 dataset shows the opposite pattern; both models perform accordingly on the MMLU and Vicuna benchmarks. This underscores the need not only for improved benchmarks and evaluation metrics but also for careful consideration of what exactly is being assessed. Are we aiming to develop models that excel in knowledge typical of high school and college settings, or do we prioritize proficiency in chatbot conversational abilities? Or perhaps another focus altogether? Since it is generally simpler to evaluate against existing benchmarks rather than creating new ones, certain benchmarks might inadvertently guide the community in a particular direction. It is crucial that, as a community, we ensure benchmarks truly reflect what we value.

While we offer a comprehensive evaluation of general chatbot performance, another limitation lies in the relatively narrow responsible AI assessment of {Guanaco}. We compare the likelihood of {Guanaco}-65B generating socially biased token sequences to other models in Table~\ref{tbl:crows}, observing that the average bias score for {Guanaco}-65B is substantially lower than that of other raw pretrained models. This suggests that finetuning on the OASST1 dataset helps mitigate the biases present in the LLaMA base model. Although these findings are promising, it remains uncertain how well {Guanaco} performs when tested against other types of biases. We defer further investigation of bias in {Guanaco} and similar chatbots to future research.

Another limitation is that we did not examine the effects of different bit-precisions, such as 3-bit base models, nor did we explore alternative adapter methods. Beyond LoRA, there exists a broad range of Parameter Efficient FineTuning (PEFT) approaches that have demonstrated effectiveness. However, it is unclear whether these methods scale effectively to very large models. We selected LoRA due to its well-established robustness, though other adapters may achieve superior results. Since finetuning after quantization appears to recover much of the information lost during quantization, this could enable even more aggressive quantization strategies. For example, 3-bit GPTQ quantization of the base model combined with LoRA finetuning might yield performance comparable to full 16-bit finetuning.

\section{Broader Impacts}

Our \textsc{QLoRA} finetuning approach is the first to enable finetuning of 33B parameter models on a single consumer GPU and 65B parameter models on a single professional GPU without sacrificing performance relative to full finetuning baselines. We have shown that our best 33B model trained on the Open Assistant dataset can match ChatGPT’s performance on the Vicuna benchmark. Since instruction finetuning is a crucial step to convert raw pretrained LLMs into ChatGPT-like chatbots, we believe our method will democratize finetuning, especially benefiting researchers with limited resources—a significant advancement for making cutting-edge NLP technology more accessible. \textsc{QLoRA} serves as a leveling mechanism that helps narrow the resource gap between large corporations and smaller teams equipped with consumer GPUs.

Another possible area of influence is the deployment on mobile devices. We propose that our \textsc{QLoRA} approach could achieve the significant breakthrough of facilitating the fine-tuning of LLMs on smartphones and other resource-constrained environments. Although 7B models have previously been demonstrated to run on phones, \textsc{QLoRA} is the first technique that permits their fine-tuning. We estimate that using an iPhone 12 Plus, \textsc{QLoRA} could fine-tune 3 million tokens overnight while the device is charging. Although fine-tuned 7B models do not match the performance of ChatGPT, we believe their quality suffices to support new applications that were previously unfeasible due to privacy concerns or LLM performance limitations. \textsc{QLoRA} facilitates privacy-conscious usage of LLMs, empowering users to control their own data and models, while simultaneously simplifying the deployment of LLMs.

Nonetheless, fine-tuning is a technology with dual-use potential and can be exploited for harmful purposes. The widespread adoption of LLMs carries well-known risks \citep{bommasani2021opportunities,bender2021dangers}, but we contend that democratizing access to this increasingly pervasive technology will foster better, more independent scrutiny compared to concentrating LLM capabilities within large corporations that do not provide models or source code for evaluation.

Overall, we anticipate that \textsc{QLoRA} will have a generally positive effect by making the fine-tuning of high-quality LLMs far more accessible and easier to perform.

\end{document}